\documentclass[11pt]{article}

\usepackage{sectsty}
\usepackage{graphicx}
\usepackage{outlines}
\usepackage{xcolor}
\usepackage{enumitem}
\setlist[enumerate,1]{label=\Roman*)}
\setlist[enumerate,2]{label=\arabic*)}
\setlist[enumerate,3]{label=\roman*)}

\usepackage{amssymb}
\newcommand\syllogism[3][]{%
  \begin{center}
  \def\tmp{#1}%
  \ifx\tmp\empty\else(#1)\quad\fi
  \begin{tabular}{@{}l@{}}
    #2\\\hline#3\quad$\therefore$
  \end{tabular}
  \end{center}
}

\usepackage{amsfonts,amsmath}
\usepackage{physics}
\usepackage{dsfont}

% Margins
\topmargin=-0.45in
\evensidemargin=0in
\oddsidemargin=0in
\textwidth=6.5in
\textheight=9.0in
\headsep=0.25in

\title{ R1 Outline Idea}
%\author{ Author }
\date{\today}

\begin{document}
\maketitle	
\pagebreak

% Optional TOC
% \tableofcontents
% \pagebreak

%--Paper--

% \section{Current Paper Outline}

% \begin{outline}[enumerate]
%     \1 Introduction

%         \syllogism[S1]{
%           Business researchers use explainers to make definitive claims about data. \\
%           They think that SHAP/LIME values \textbf{describe} the data. \\
%           They think that SHAP/LIME values count as evidence upon which \textbf{prescriptive} claims can be made. \\
%           They do not suppose a mathematical definition when using the term ``importance." \\
%           Previous research shows explainers cannot even explain $M$. \\
%           We showed how and why explainers fail to give descriptive/prescriptive insights into $G$.
%         }{
%           Business researchers ought only to use explainers to \textbf{generate hypotheses} about the data (mention emerging field of research, including those using LLM for hypotheses) cite the QJE paper on hypotheses proposal from image gradient. \\
%         }

%     \1 Related Work
%             need to add discussion on how econometrics methods are used to identify relationships between features and the target variable in data
%     \1 Use of SHAP and LIME in Business Research
%         \2 Uses with \emph{Descriptive} View of Feature Importance
%             \3 Infer that $X_j$ more/less important than $X_k$ in DGP 
%         \2 Uses with \emph{Prescriptive} View of Feature Importance
%             \3  Causal Claims like Counterfactuals: ``If Stakeholders \emph{were} to perform intervention $X\mapsto X'$ then it \emph{would} be the case that $G(X')= Y^*$
        

%     \1 Explanation Pipeline
%         \2 Introduce Stage 1 \& Stage 2
        
%         \2 Introduce synthetic data generation

%         \2 Applying Pipeline

%         \2 Pervasiveness of Misalignment

%     \1 Obtaining Descriptive and Prescriptive Insights into a Linear Ground Truth Model
    
%         \2 Metrics for Linear Ground Truth
%             \3 Descriptive: Directionality
%             \3 Descriptive: Concordance
%             \3 Descriptive: Relevance

%         \2 sources of discrepancy (why SHAP and LIME cannot explain the data?)
% Stage 1: SHAP and LIME quantify feature importance within the learned model, not within the underlying data or data-generating process, and the model's decision-making process does not necessarily align with the data generating process.
% (we want to show that there exists so many mappings from x to y, and we have no idea which is the true model's)

% Dataset Properties (number of features, feature correlations)

% Predictive Power of $M$ - even you have perfect accuracy, you may not completely align because of the rashomon effect
% (ask gpt if there's)

% Stage 2: The ways in which researchers often use SHAP and LIME—such as inferring which features are important in the data or reasoning about how modifying a feature would affect the outcome—do not align with the formal definitions or intended interpretations of these methods.


%     \1 Obtaining Descriptive and Prescriptive Insights into Nonlinear Ground Truth Models

%         \2 General Characterizations of Ground Truth Feature Importance
%             \3 LIME/SHAP val of $G(x)$ (\textcolor{red}{done})
%             % \3 Local: $\nabla G\mid_x$ 
%             % https://arxiv.org/pdf/2201.08837
%             \3 Global: Average Marginal Effect (AME) of $G$: the conditional expectation of Y given X, focusing on the *average* marginal effect across the population.
%             % \3 Global: Marginal Effect of $G$ at Mean Vector (MMR)
%             % \3 Global: Marginal Effect of $G$ at Representative Value (MER)

%         \2 Dataset Construction

%         \2 Metrics for Nonlinear Ground Truth
%             \3 Descriptive: Rank Concordance (Local + Global)
%             \3 \textbf{Descriptive: Directionality using AME/MMR/MER (Local + Global)}
%             \3 Prescriptive: Counterfactual Validity (cite the MISQ paper)
            

%         \2  Experiments

%     \1 Robustness Checks

%     \1 Discussion and Conclusion
%         \2 Main Findings

%         \2 On the Use of Post Hoc Explanations
%                 talk about using it for hypotheses generation and cite emergent work on it
%         \2 General Concerns


%         \2 Final Remarks
% \end{outline}

% \pagebreak

% \section{Argument Structures}

% \subsection{v1}

%         \syllogism[S1]{
%           Business researchers use explainers to make definitive claims about data. \\
%           They think that SHAP/LIME values \textbf{describe} the data. \\
%           They think that SHAP/LIME values count as evidence upon which \textbf{prescriptive} claims can be made. \\
%           They do not suppose a mathematical definition when using the term ``importance." \\
%           Previous research shows explainers cannot even explain $M$. \\
%           We showed how and why explainers fail to give descriptive/prescriptive insights into $G$.
%         }{
%           Business researchers ought only to use explainers to \textbf{generate hypotheses} about the data (mention emerging field of research, including those using LLM for hypotheses) cite the QJE paper on hypotheses proposal from image gradient. \\
%         }

%     \subsection{v2}

%         \textbf{Reduce emphasis on CORRECTNESS of explanations}

%         \syllogism[S2]{
%           Business researchers use explainers to derive a theoretical model $T'$ of a social phenomenon. \\
%           They want $T'$ to have fewer internal theoretical inconsistencies than their previous theory $T$ had. \\
%           They want to be able to use $T'$ to solve more empirical problems than they could using their previous theory $T$. \\ 
%           % They think that SHAP/LIME values \textbf{describe} the data. \\
%           % They think that SHAP/LIME values count as evidence upon which \textbf{prescriptive} claims can be made. \\
%           % They do not suppose a mathematical definition when using the term ``importance." \\
%           Previous research shows explainers cannot even explain $M$. \\
%           We showed how and why explainers fail to give descriptive/prescriptive insights into $G$.
%         }{
%           Business researchers ought only to use explainers to \textbf{generate hypotheses} about the data, \\\emph{must} check those hypotheses for contradictions, and \emph{should} test the hypotheses experimentally.   \\
%         }


% ===





\section{Introduction}
\section{Related Work}
\section{Current Uses and Interpretations of SHAP and LIME in Business Research}
\subsection{Prevalence of Data-Level Inference}
summary statistics of the misuses, especially in top journals (Lu)
\subsection{The Explanation Pipeline in Business Research}
Talk about the pipeline, use one linear ground truth data as an example to demonstrate the explanations of SHAP and LIME

show the SHAP plot (beeswarm plot) and coefficient plot for LIME
\subsection{Types of Interpretations Drawn from SHAP and LIME}
\subsubsection{SHAP Explanations}
Claim 1: Larger shapley values indicate larger Y

Claim 2: features with higher mean absolute SHAP values are more important

\subsubsection{LIME explanations}
LIME explanations approximate local marginal effects

\section{Do Post-hoc Explanations Align with Data?}
\subsection*{SHAP estimate marginal contribution}
Shapley value represents the contribution of a feature to the dependent variable
\paragraph{Claim 1: large SHAP values indicate larger (or more positive) $Y$ - \emph{Direction of Contribution}}

evaluation:  Spearman $\rho$ and Kendall $\tau$ between shapley value ranking and $G(x_i)$ - do larger shapley values co-occur with larger outcomes? (this is the part that is undone)

    Define the perturbation function $p^j_\delta(\mathbf{x}) := [x_1,x_2,\ldots,x_j+\delta,\ldots,x_d]$, and we can generate many different perturbed data points, such as $p^j_2\delta(\mathbf{x}), p^j_3\delta(\mathbf{x}), \cdots p^j_{-1\delta(\mathbf{x})}, p^j_{-2\delta(\mathbf{x})}, \cdots$.
    For each of these ``counterfactual" or "perturbed" data points, we compute the shapley value for feature j, so we have $e_{j,\delta}, e_{j,2\delta}, \cdots$. We also have the ground truth $Y$ for each data point, $G(p^j_\delta(\mathbf{x})), G(p^j_2\delta(\mathbf{x}))$, etc. We compare the ranking based on $e_{j,\delta}$ with the ranking based on $G(p^j_\delta(\mathbf{x}))$, for a fixed $j$. The idea is, as the above interpretation claims ``if you see a larger shapley value (given everything else fixed), you should see a larger $Y$. If we rank the data points based on $e_{j}$ in an ascending order, we expect to see $Y$ also increases or the probability of positive increases.
    
    % Define $\rho,\ldots,\rho_d$ so that $\rho_j = \text{corr}(\{e^{\text{SHAP}}_j(\mathcal{M}(p_{k\delta}^j(\mathbf{x}))\}_{k=1}^{10}, \{(\sigma\circ \mathcal{G})(p_{k\delta}^j(\mathbf{x})) - (\sigma \circ \mathcal{G})(\mathbf{x})\}_{k=1}^{10})$

    % Thus $\rho_j$ can be used to check how similar the effect of increasing a $X_j$ while holding the other features constant are on SHAP values and on increases to $\mathds{P}[Y=1]$.


\paragraph{Claim 2: SHAP identifies important features based on mean absolute SHAP value}
evaluation: Spearman $\rho$ and Kendall $\tau$ between ranking based on mean absolute SHAP values and true shapley from G
experiments done by Lu Feng 


\subsection*{LIME estimates marginal effect}

evaluation: directionality, relative importance (old results from 1st submission and need to do new ones for non-linear)


\subsection{Data Simulation}
\subsubsection{Results of SHAP evaluation}
\subsubsection{Results of LIME evaluation}
\subsection{Results}




\section{why it happened}
M is not G (we can just focus on SHAP, put LIME in the appendix)

\subsection{the number of features}

\subsection{correlation of features}

\subsection{the M's accuracy }
    1) M is not perfect: alignment gets better when accuracy increases (experiments)
    2) even M has high accuracy, the Rashomon set theory (discussion) -> show evidence of m highly accurate models disagree in explanations

\subsection{Explainer Properties}


regression analyses

\section{Robustness Check}
1. Average explanation (e.g., 10 M, each with a shap explanation)
2. Average the model first and then explain the ensemble
3. for LIME, average LIME with different seeds

take-away: talk about using agreement among the models as a sign for not trusting or generalizing the explanations from the model to the data. 
\section{Discussion}


call for future work on "confidence interval" or "p-value" for model explanations

suggestions on diagnoses (how do we know we should not trust the explanation?)

advocate for using post-hoc explanations for hypotheses generation rather than evidence



===
What we need from Lu:

Data generation code
were models saved? were shapley values saved?
\end{document}